\chapter{Conception de l'organisation de l'ingénierie}\label{ch:fm:engineering-organisation-design}

En avril 1968, Melvin Conway publie l'observation selon laquelle «~les organisations qui conçoivent des systèmes \dots\ sont contraintes de produire des conceptions qui sont
des copies des structures de communication de ces organisations~». Une société de trading qui répartit ses ingénieurs par classe d'actifs finit avec trois
passerelles d'ordres, trois contrôles de risque et trois gestionnaires de données de marché, qu'elle le veuille ou non, car trois équipes qui se parlent peu construisent trois
exemplaires de chaque chose. L'organigramme est le premier brouillon de l'architecture.

\section{Équipes intégrées contre équipes centrales}

\begin{definition}[Équipe intégrée, équipe centrale]\label{def:fm:engineering-organisation-design:teams}
Une \emph{équipe intégrée}\index{equipe integree@équipe intégrée} (embedded team) est rattachée à un métier (un desk ou un groupe de stratégies), en rend compte à celui-ci et construit tout
ce dont ce métier a besoin. Une \emph{équipe centrale}\index{equipe centrale@équipe centrale} (central team) sert plusieurs métiers avec une même fonction, telle que les passerelles d'ordres, les données de marché ou les contrôles de risque, et
rend compte à la technologie.
\end{definition}

Les deux conceptions échangent les mêmes choses. Les \omterm{def:fm:engineering-organisation-design:teams}{équipes intégrées} parlent à leurs traders tous les jours, évoluent vite et possèdent l'ensemble de leur pile~; elles dupliquent
ce que d'autres desks construisent aussi, et chaque copie demande sa propre maintenance et sa propre astreinte. Les \omterm{def:fm:engineering-organisation-design:teams}{équipes centrales} construisent une chose bien pour tout le monde et peuvent
consolider les doublons~; chaque modification dont un desk a besoin de leur part franchit une frontière d'équipe. Les \omterm{def:fm:technology-strategy:teams}{équipes plateforme} et produit du chapitre~19 sont les noms modernes des
deux pôles, et \emph{Team Topologies} de Skelton et Pais (2019) en fournit le vocabulaire~; MacCormack, Rusnak et Baldwin (2012) ont testé le
«~miroir~» de Conway, c'est-à-dire la reproduction de l'organisation dans l'architecture, de façon empirique.

\begin{definition}[Coût de coordination]\label{def:fm:engineering-organisation-design:coord}
Le \emph{coût de coordination}\index{coût de coordination} (coordination cost) d'une conception organisationnelle est le travail nécessaire pour garder les équipes en phase à travers les dépendances
entre les services qu'elles possèdent~; le chapitre le mesure comme le nombre de modifications par semaine portant sur des services dont dépendent les services d'autres équipes, lesquelles
doivent connaître, tester et intégrer à leur planning.
\end{definition}

\section{Le développeur quant et les autres rôles hybrides}

\begin{definition}[Développeur quant]\label{def:fm:engineering-organisation-design:qd}
Un \emph{développeur quant}\index{developpeur quant@développeur quant} (quant developer) est un ingénieur qui implémente et fait tourner des modèles de trading~: il transforme le prototype d'un chercheur en
code de production, répond de sa performance et de sa justesse, et connaît assez le modèle pour le modifier sans risque.
\end{definition}

Les rôles hybrides existent parce que la transmission entre la recherche et la production est l'endroit où les modèles cassent. Le livre~12 a décrit l'ingénieur en apprentissage automatique (ML engineer)
et la spécification de transmission (handoff) entre un propriétaire du modèle et la plateforme~; le \omterm{def:fm:engineering-organisation-design:qd}{développeur quant} est le même rôle pour tout type de modèle. Dans une conception
intégrée, il se trouve au sein du desk~; dans une conception centrale, il se situe entre le desk et la plateforme, et c'est sur lui que retombe le \omterm{def:fm:engineering-organisation-design:coord}{coût de coordination}.

\section{Propriété~: chaque service a un responsable}

\begin{definition}[Responsable de service, facteur bus]\label{def:fm:engineering-organisation-design:owner}
Un \emph{responsable de service}\index{responsable de service} (service owner) est l'équipe désignée, avec un responsable nommé, qui répond de la justesse d'un service, de ses modifications, de son astreinte et
de son retrait. Le \emph{facteur bus}\index{facteur bus} (bus factor) d'un service est le nombre de personnes qui peuvent l'exploiter et le modifier~; le facteur bus d'une
organisation est le plus petit de ces facteurs sur l'ensemble de ses services~: le nombre de départs qui laisserait un service sans personne qui le comprenne.
\end{definition}

La propriété transforme la conception en données (voir le \cref{lst:fm:engineering-organisation-design:measures})~: un tableau des services avec leur responsable, leurs
dépendances, la fréquence de chaque modification et celle de leurs appels. Il en découle trois mesures~: le \omterm{def:fm:engineering-organisation-design:coord}{coût de coordination}, la charge d'astreinte
par ingénieur et le \omterm{def:fm:engineering-organisation-design:owner}{facteur bus}.

\omcode{code/firm/teamtopo/firm_teamtopo.py}{34}{61}{Les trois mesures d'une conception organisationnelle~: dépendances inter-équipes pondérées par
la fréquence des modifications, appels par ingénieur par équipe, et facteur bus.}\label{lst:fm:engineering-organisation-design:measures}

\begin{proposition}[Les conceptions déplacent la charge, les services la créent]\label{prop:fm:engineering-organisation-design:load}
Pour un ensemble fixé de services et d'ingénieurs, la charge d'astreinte moyenne par ingénieur est la même dans toutes les conceptions~: le total des appels divisé par le nombre d'ingénieurs. Une
conception ne change que la répartition de la charge entre les équipes, le \omterm{def:fm:engineering-organisation-design:coord}{coût de coordination} et le \omterm{def:fm:engineering-organisation-design:owner}{facteur bus}. Seule la suppression ou la fusion de services, ou
la réduction de leurs appels, abaisse la moyenne.
\end{proposition}

\begin{proof}
Les appels de chaque service sont comptés une seule fois, dans l'équipe de son responsable, quel qu'il soit~; en sommant sur les équipes on obtient le total des appels, et en divisant par les
ingénieurs, qui appartiennent tous à une équipe, donne la moyenne.
\end{proof}

\section{Tutoriel~: trois passerelles d'ordres}\label{tut:fm:engineering-organisation-design:tut}

\begin{tutorial}
\textbf{Objectif.} Comparer une conception intégrée et une conception centrale pour vingt services et quinze ingénieurs, puis consolider, puis déplacer un ingénieur et
observer ce qui casse. \textbf{État final~:} le \cref{tab:fm:engineering-organisation-design:designs} et le graphique de charge
(la \cref{fig:fm:engineering-organisation-design:pages}).
\begin{tutsteps}
\item \textbf{Les services.} Trois desks (actions, contrats à terme, options), chacun avec une stratégie, une passerelle d'ordres, un contrôle de risque, un gestionnaire de données de marché et
une bibliothèque de pricing~; cinq services partagés (référentiel des titres, post-marché, supervision, déploiement, plateforme de recherche)~; 33 dépendances~; 34.5 modifications et
6.2 appels par semaine au total (\texttt{fm\_teams}, illustratif).
\item \textbf{Intégrée.} L'équipe de quatre personnes de chaque desk possède ses cinq services~; une équipe d'infrastructure de trois personnes possède les services partagés.
\item \textbf{Centrale.} Les équipes de desk de deux personnes possèdent les stratégies et le pricing~; des équipes de deux personnes possèdent les trois passerelles, les trois contrôles de risque et les trois
gestionnaires de données de marché~; l'infrastructure reste la même.
\item \textbf{Consolidée.} La suite naturelle de la conception centrale~: une passerelle, un contrôle de risque et un gestionnaire de données de marché, possédés par une \omterm{def:fm:technology-strategy:teams}{équipe plateforme} de trading de
trois personnes~; des équipes de desk de trois personnes~; quatorze services.
\item \textbf{Le déplacement.} Un ingénieur quitte un desk (conception intégrée) ou l'équipe des passerelles (conception centrale) pour rejoindre l'infrastructure.
\end{tutsteps}
\end{tutorial}

\begin{table}[ht]
\centering\small
\begin{tabular}{>{\raggedright\arraybackslash}p{2.3cm}>{\centering\arraybackslash}p{1.8cm}>{\centering\arraybackslash}p{1.8cm}>{\centering\arraybackslash}p{1.8cm}>{\centering\arraybackslash}p{1.8cm}>{\centering\arraybackslash}p{1.8cm}>{\centering\arraybackslash}p{1.8cm}}
\toprule
 & & & inter-équipes & coordi\-nation & appels max. par & facteur \\
conception & équipes & services & arêtes & (modif. par semaine) & ingénieur et semaine & bus \\
\midrule
intégrée & 4 & 20 & 18 & 22.5 & 0.43 & 3 \\
centrale & 7 & 20 & 30 & 33.0 & 0.90 & 2 \\
centrale, consolidée & 5 & 14 & 19 & 29.0 & 0.67 & 3 \\
\bottomrule
\end{tabular}
\caption{Trois conceptions pour quinze ingénieurs. Moyenne des appels par ingénieur et par semaine~: 0.41 pour les deux conceptions des vingt services, 0.29 après
consolidation. Données~: \texttt{fm\_teams.compare\_all}.}\label{tab:fm:engineering-organisation-design:designs}
\end{table}

La conception intégrée offre la coordination la plus faible (22.5 modifications par semaine franchissant les équipes, sur 18 arêtes), la charge d'astreinte la plus régulière (0.43 appel par
ingénieur et par semaine au maximum) et un \omterm{def:fm:engineering-organisation-design:owner}{facteur bus} de 3. Son coût est invisible dans le tableau~: elle fait tourner trois exemplaires de tout. Centraliser les mêmes
vingt services est le pire des deux mondes~: 30 arêtes inter-équipes, 33 modifications par semaine à coordonner, l'équipe des gestionnaires de données de marché portant 0.90 appel par ingénieur et par semaine, et
un \omterm{def:fm:engineering-organisation-design:owner}{facteur bus} de 2, puisque chaque fonction repose désormais sur deux personnes. La charge moyenne ne bouge pas
(voir la \cref{prop:fm:engineering-organisation-design:load})~: réorganiser seul redistribue le travail et ajoute des frontières.

La consolidation est ce à quoi sert la conception centrale. Fusionner les trois passerelles, contrôles de risque et gestionnaires de données de marché laisse quatorze services et 4.3
appels par semaine, soit 31~\% de moins~; la charge moyenne descend à 0.29 appel par ingénieur, le \omterm{def:fm:engineering-organisation-design:owner}{facteur bus} revient à 3 avec des équipes de trois personnes, et les desks gagnent chacun un
ingénieur. La coordination reste plus élevée que dans la conception intégrée (29.0 modifications par semaine), car chaque desk dépend désormais d'une seule passerelle partagée~:
c'est le prix d'en avoir une.

\begin{omfigure}
\begin{tikzpicture}
\begin{axis}[width=10.5cm,height=5.2cm,ybar,bar width=9pt,ylabel={\small nombre par semaine / arêtes},ymin=0,ymax=36,xtick={0,1,2},
  xticklabels={intégrée,centrale,{centrale, consolidée}},xmin=-0.6,xmax=2.6,
  tick label style={font=\footnotesize},grid=major,grid style={gray!20},table/col sep=comma,
  legend cell align=left,legend style={font=\scriptsize,at={(0.5,-0.25)},anchor=north,/tikz/every even column/.append style={column sep=8pt}},legend columns=2]
\addplot[fill=omDef!55,draw=omDef] table[x=pos,y=coordination]{figdata/desk/21-engineering-organisation-design/designs.csv};
\addplot[fill=omProp!55,draw=omProp] table[x=pos,y=cross_edges]{figdata/desk/21-engineering-organisation-design/designs.csv};
\legend{coordination (modifications par semaine),arêtes inter-équipes}
\end{axis}
\end{tikzpicture}

\omcaption{Coordination selon les trois conceptions~: le nombre de dépendances qui franchissent une frontière d'équipe, et les modifications par semaine qui les franchissent.
Centraliser sans consolider ajoute des frontières~; consolider en supprime à nouveau une partie. Données~: \texttt{fm\_teams.compare\_all}.}\label{fig:fm:engineering-organisation-design:coord}
\end{omfigure}

\section{L'astreinte et son coût}

La charge d'astreinte se mesure en appels, mais elle se paie en heures. Le livre de Google sur la fiabilité des sites plafonne l'astreinte au quart du temps d'un ingénieur et
constate qu'un incident, avec son analyse de cause racine, sa remédiation et son suivi, prend environ six heures~; deux de ces incidents au plus tiennent dans un
poste de douze heures (voir l'\cref{dat:fm:engineering-organisation-design:sre}). Avec les chiffres du chapitre, l'équipe des données de marché de la conception centrale consacre
$0.9\times6=5.4$ heures par semaine et par ingénieur aux incidents, soit 13.5~\% d'une semaine de quarante heures~; après le départ d'un ingénieur de l'équipe des passerelles, son seul ingénieur restant
porte 1.5 appels par semaine, soit neuf heures, 22.5~\% de la semaine et proche du plafond (voir la \cref{fig:fm:engineering-organisation-design:pages}).

\begin{dated}{2026-09}{Un budget d'astreinte publié}\label{dat:fm:engineering-organisation-design:sre}
Google's \emph{Site Reliability Engineering} (2016), chapitre «~Being On-Call~»~: au moins 50~\% du temps des ingénieurs de fiabilité des sites va à
l'ingénierie et au plus 25~\% à l'astreinte~; la prise en charge d'un incident prend environ six heures en moyenne, d'où un maximum de deux incidents par
poste de douze heures. Le livre~15 applique ces budgets à la rotation d'astreinte d'une plateforme de trading.
\end{dated}

\begin{omfigure}
\begin{tikzpicture}
\begin{axis}[width=10.5cm,height=5.4cm,xbar,bar width=6pt,xmin=0,xmax=1.1,xlabel={\small appels par ingénieur et par semaine},
  symbolic y coords={embedded: desk,embedded: infrastructure,central: desk,central: gateways,central: risk,central: market data,central: infrastructure,consolidated: desk,consolidated: platform,consolidated: infrastructure},
  ytick={embedded: desk,embedded: infrastructure,central: desk,central: gateways,central: risk,central: market data,central: infrastructure,consolidated: desk,consolidated: platform,consolidated: infrastructure},yticklabels={{intégrée~: desk},{intégrée~: infrastructure},{centrale~: desk},{centrale~: passerelles},{centrale~: risque},{centrale~: données de marché},{centrale~: infrastructure},{consolidée~: desk},{consolidée~: plateforme},{consolidée~: infrastructure}},
  y dir=reverse,tick label style={font=\footnotesize},grid=major,grid style={gray!20},enlarge y limits=0.06,nodes near coords,
  nodes near coords style={font=\scriptsize,/pgf/number format/.cd,fixed,precision=2},point meta=x,table/col sep=comma]
\addplot[fill=omDef!55,draw=omDef] table[y=label,x=pages]{figdata/desk/21-engineering-organisation-design/pages.csv};
\end{axis}
\end{tikzpicture}

\omcaption{Charge d'astreinte par ingénieur et par équipe selon les trois conceptions (les trois équipes de desk portent la même charge dans chaque conception). La conception centrale
concentre les appels sur les petites équipes fonctionnelles. Données~: \texttt{firm.teamtopo.pages}.}\label{fig:fm:engineering-organisation-design:pages}
\end{omfigure}

Le transfert montre le \omterm{def:fm:engineering-organisation-design:owner}{facteur bus} à l'œuvre. Dans la conception intégrée, transférer un ingénieur hors du desk actions laisse ses cinq services avec trois
personnes et porte la charge de ce desk à 0.57 appel par semaine chacun~; aucun service n'est orphelin. Dans la conception centrale, transférer un ingénieur hors de l'équipe des passerelles
laisse les trois passerelles avec une seule personne chacune~: le \omterm{def:fm:engineering-organisation-design:owner}{facteur bus} tombe à 1, et une seule démission laisserait la société incapable de modifier la façon dont elle envoie les
ordres. Dans la conception consolidée, le même transfert laisse la plateforme de trading avec deux personnes et un \omterm{def:fm:engineering-organisation-design:owner}{facteur bus} de 2.

\begin{omfigure}
\begin{tikzpicture}[font=\footnotesize,>=stealth,t/.style={draw,rounded corners,fill=omDef!15,minimum width=1.6cm,align=center},
  s/.style={draw,rounded corners,fill=omThm!15,minimum width=1.5cm,align=center}]
\node[font=\small] at (-4.2,2.5) {intégrée};
\node[t] (e) at (-6,1.5) {actions}; \node[t] (f) at (-4.2,1.5) {future}; \node[t] (o) at (-2.4,1.5) {options};
\node[s] (ge) at (-6,0) {passerelle 1}; \node[s] (gf) at (-4.2,0) {passerelle 2}; \node[s] (go) at (-2.4,0) {passerelle 3};
\draw[->] (e) -- (ge); \draw[->] (f) -- (gf); \draw[->] (o) -- (go);
\node[font=\small] at (3.6,2.5) {centrale, consolidée};
\node[t] (e2) at (1.8,1.5) {actions}; \node[t] (f2) at (3.6,1.5) {future}; \node[t] (o2) at (5.4,1.5) {options};
\node[t,fill=omProp!15] (p) at (3.6,0) {équipe plateforme de trading}; \node[s] (g) at (3.6,-1.3) {une passerelle};
\draw[->,dashed] (e2) -- (p); \draw[->,dashed] (f2) -- (p); \draw[->,dashed] (o2) -- (p); \draw[->] (p) -- (g);
\draw[gray] (-0.3,-1.8) -- (-0.3,2.8);
\end{tikzpicture}

\omcaption{La loi de Conway dans une société de trading~: trois équipes de desk construisent trois passerelles (à gauche)~; une \omterm{def:fm:technology-strategy:teams}{équipe plateforme} en construit une seule, dont
dépendent les trois desks (en pointillés, à droite). La seconde conception compte moins de services et davantage de dépendances inter-équipes. Schéma de principe.}\label{fig:fm:engineering-organisation-design:conway}
\end{omfigure}

\section{Mesurer une organisation d'ingénierie}

Les trois mesures sont peu coûteuses à tenir dès lors que la propriété des services est une donnée~: le \omterm{def:fm:engineering-organisation-design:coord}{coût de coordination} se calcule à partir du graphe de dépendances et des journaux de modifications, la charge d'astreinte à partir du
système d'alertes, le \omterm{def:fm:engineering-organisation-design:owner}{facteur bus} à partir de ceux qui ont modifié et exploité chaque service au cours de l'année écoulée. Le livre de Forsgren, Humble et Kim, \emph{Accelerate} (2018),
propose des mesures de livraison (la fréquence des déploiements, le délai qu'une modification met pour atteindre la production, la fréquence des échecs, la rapidité
de rétablissement du service) qui les complètent.

\begin{method}[Concevoir l'organisation d'ingénierie]\label{met:fm:engineering-organisation-design:design}
\begin{enumerate}
\item Tenir le catalogue des services sous forme de données~: propriétaire, dépendances, fréquence des modifications, taux d'appels, qui peut l'exploiter.
\item Dessiner les équipes autour des échanges qu'exige le travail (critère de Conway)~: intégrer ce qui évolue avec le desk, centraliser ce dont plusieurs desks
ont besoin sous la même forme, et consolider ce qui est dupliqué.
\item Mesurer la coordination, la charge d'astreinte par rapport à un budget, et le \omterm{def:fm:engineering-organisation-design:owner}{facteur bus}, avant et après chaque réorganisation.
\item Maintenir le \omterm{def:fm:engineering-organisation-design:owner}{facteur bus} de chaque service à trois au moins~; faire binômer en priorité les nouveaux arrivants avec les services les plus fragiles.
\item Revoir la conception lorsque l'on ajoute un desk ou une plateforme de négociation~: chacun est un nouveau service ou une nouvelle dépendance.
\end{enumerate}
\end{method}

\section{Construction~: l'organisation comme graphe}\label{bld:fm:engineering-organisation-design:teamtopo}

\begin{build}
\bfield{Rôle} Services, dépendances, taux de modification et d'incidents et affectations aux équipes, sous forme de données, avec les mesures d'une conception.
\bfield{Interface} \texttt{firm.teamtopo}~: \texttt{Service}, \texttt{Design}, \texttt{cross\_edges}, \texttt{coordination}, \texttt{pages},
\texttt{bus\_factor}, \texttt{move}, \texttt{report}.
\bfield{Règles} Chaque service a exactement une équipe propriétaire~; le \omterm{def:fm:engineering-organisation-design:owner}{facteur bus} d'un service compte son équipe et les suppléants nommés~; la charge moyenne est invariante
lors d'une réaffectation.
\bfield{Tests d'acceptation} \texttt{code/firm/teamtopo/tests/}~: arêtes inter-équipes et coordination sur un graphe à trois services~; appels par équipe~; le \omterm{def:fm:engineering-organisation-design:owner}{facteur bus}
avant et après un transfert.
\bfield{Pour aller plus loin} Connaissances mesurées à partir de l'historique des commits et des déploiements~; le coût de la réorganisation elle-même~; optimisation de l'affectation pour une
somme pondérée des trois mesures.
\end{build}

\begin{omsources}
\item M.~E.~Conway, «~How do committees invent?~», \emph{Datamation}, avril 1968.
\item A.~MacCormack, C.~Baldwin et J.~Rusnak, \emph{Research Policy} 41(8), 2012~; M.~Skelton et M.~Pais, \emph{Team Topologies}, 2019.
\item B.~Beyer et al. (dir.), \emph{Site Reliability Engineering}, 2016~; N.~Forsgren, J.~Humble et G.~Kim, \emph{Accelerate}, 2018.
\end{omsources}

\section{Exercices}

\begin{exercise}[$\star$]\label{exo:fm:engineering-organisation-design:1}
Énoncer la thèse de Conway et son critère pour organiser un effort de conception.
\end{exercise}

\begin{exercise}[$\star$]\label{exo:fm:engineering-organisation-design:2}
Quelle est la charge moyenne d'astreinte par ingénieur dans chaque conception, et pourquoi est-elle la même pour les conceptions intégrée et centrale~?
\end{exercise}

\begin{exercise}[$\star$]\label{exo:fm:engineering-organisation-design:3}
Combien d'heures par semaine l'\omterm{def:fm:engineering-organisation-design:teams}{équipe centrale} de données de marché consacre-t-elle aux incidents, à raison de six heures par incident, et quelle part du plafond d'astreinte cela représente-t-il~?
\end{exercise}

\begin{exercise}[$\star\star$]\label{exo:fm:engineering-organisation-design:4}
Pourquoi centraliser les mêmes vingt services fait-il passer la coordination de 22.5 à 33.0 changements par semaine~?
\end{exercise}

\begin{exercise}[$\star\star$]\label{exo:fm:engineering-organisation-design:5}
Que devient le \omterm{def:fm:engineering-organisation-design:owner}{facteur bus} dans chaque conception lorsqu'un ingénieur change d'équipe, et quels services sont exposés~?
\end{exercise}

\begin{exercise}[$\star\star$]\label{exo:fm:engineering-organisation-design:6}
Où les \omterm{def:fm:engineering-organisation-design:qd}{développeurs quant} doivent-ils se situer dans chaque conception, et pourquoi~?
\end{exercise}

\begin{exercise}[$\star\star\star$]\label{exo:fm:engineering-organisation-design:7}
\emph{Programmation.} Dans la conception consolidée, transférer un ingénieur de la plateforme de trading vers l'infrastructure. Quels sont le \omterm{def:fm:engineering-organisation-design:owner}{facteur bus} et la
charge la plus élevée~?
\end{exercise}

\begin{exercise}[$\star\star\star$]\label{exo:fm:engineering-organisation-design:8}
\emph{Trouver la faille.} «~Centraliser nos ingénieurs réduira notre charge d'astreinte.~»
\end{exercise}

\section{Problème~: trois passerelles d'ordres}

\begin{problem}[{Problème du week-end --- trois passerelles d'ordres}]\label{pb:fm:engineering-organisation-design:1}
Un nouveau directeur technique hérite de trois desks, chacun avec ses propres ingénieurs et trois exemplaires de tout, et doit proposer une organisation.

\medskip\noindent\textbf{Partie I --- Les principes.}
\begin{enumerate}
\item Énoncer la thèse de Conway.
\item Définir les \omterm{def:fm:engineering-organisation-design:teams}{équipes intégrées} et centrales, ainsi que leur compromis.
\item Définir un \omterm{def:fm:engineering-organisation-design:qd}{développeur quant} et expliquer pourquoi ce rôle existe.
\item Définir le \omterm{def:fm:engineering-organisation-design:coord}{coût de coordination}, le \omterm{def:fm:engineering-organisation-design:owner}{responsable de service} et le \omterm{def:fm:engineering-organisation-design:owner}{facteur bus}.
\end{enumerate}

\medskip\noindent\textbf{Partie II --- Le modèle.}
\begin{enumerate}[resume]
\item Décrire les services et leurs dépendances.
\item Énoncer et démontrer la \cref{prop:fm:engineering-organisation-design:load}.
\item Donner les trois mesures pour les conceptions intégrée et centrale.
\item Les donner pour la conception consolidée, et expliquer ce que la consolidation apporte et ce qu'elle coûte.
\end{enumerate}

\medskip\noindent\textbf{Partie III --- Stress.}
\begin{enumerate}[resume]
\item Quel est le budget d'astreinte publié, et à quel point chaque conception s'en approche-t-elle~?
\item Transférer un ingénieur dans chaque conception et donner le résultat.
\item Quels services faut-il appairer en premier, et pourquoi~?
\item Comment mesurer le \omterm{def:fm:engineering-organisation-design:owner}{facteur bus} à partir des données~?
\end{enumerate}

\medskip\noindent\textbf{Partie IV --- La proposition.}
\begin{enumerate}[resume]
\item Quels services doivent rester intégrés, lesquels doivent être centralisés, lesquels doivent être consolidés~?
\item Comment séquencer le changement~?
\item Quelles mesures de livraison ajouter~?
\item Quel est le risque d'une passerelle unique partagée, et comment le gérer~?
\item À quelle fréquence réviser la conception~?
\item Pourquoi réorganiser sans consolider constitue-t-il la pire option du modèle~?
\item Énoncer le \emph{résultat nommé}~: le \omterm{def:fm:engineering-organisation-design:coord}{coût de coordination} et les pages par ingénieur et par semaine dans les conceptions intégrée et centrale, ainsi que le \omterm{def:fm:engineering-organisation-design:owner}{facteur bus}
de chacune.
\item En deux phrases, rédiger la proposition.
\end{enumerate}
\end{problem}

\section{Questions d'entretien}

\begin{interviewq}[$\star$ \iqroles{développeur}]\label{iq:fm:engineering-organisation-design:1}
Qu'est-ce que la loi de Conway, et où l'avez-vous déjà observée~?
\end{interviewq}

\begin{interviewq}[$\star$ \iqroles{développeur, chercheur}]\label{iq:fm:engineering-organisation-design:2}
Que fait un \omterm{def:fm:engineering-organisation-design:qd}{développeur quant} que ni un chercheur ni un ingénieur logiciel ne font~?
\end{interviewq}

\begin{interviewq}[$\star\star$ \iqroles{développeur}]\label{iq:fm:engineering-organisation-design:3}
Comment calculeriez-vous le \omterm{def:fm:engineering-organisation-design:owner}{facteur bus} d'une base de code~?
\end{interviewq}

\begin{interviewq}[$\star\star$ \iqroles{développeur}]\label{iq:fm:engineering-organisation-design:4}
Votre équipe reçoit deux pages par nuit. Que faites-vous~?
\end{interviewq}

\begin{interviewq}[$\star\star$ \iqroles{trader, développeur}]\label{iq:fm:engineering-organisation-design:5}
Chaque desk de trading doit-il avoir ses propres ingénieurs~?
\end{interviewq}

\begin{interviewq}[$\star\star\star$ \iqroles{développeur, chercheur}]\label{iq:fm:engineering-organisation-design:6}
Formulez l'affectation des services aux équipes comme un problème d'optimisation.
\end{interviewq}
